\chapter{Gaussian Integrals and Convolution}\label{ch:gaussian-convolution}

\section{Finite convolution and moments}

For positive rational weights of total mass one, form the list of all
pairwise sums of sample points and multiply their weights. If
$A_K(h)=\sum_i p_i h(x_i)$, the exact finite identity is
\[
 A_{K*L}(h)=\sum_i p_i\sum_jq_jh(x_i+y_j).
\]
Finite interchange proves commutativity and associativity of these actions,
without identifying different list encodings. Means and variances satisfy
\[
 \mu_{K*L}=\mu_K+\mu_L,\qquad v_{K*L}=v_K+v_L,
 \qquad v_{K^{*n}}=nv_K.
\]
The function action $T_Kh(x)=\sum_i p_i h(x-x_i)$ satisfies
$T_{K*L}h=T_K(T_Lh)$.

\begin{theorem}[Finite convolution laws]
\label{thm:finite-convolution-laws}
The pairwise-sum construction preserves nonnegative weights and mass one.
Its actions are commutative and associative, and the moment and function
identities above hold for every supplied finite probability kernel.
\lean{ComputableAnalysis.WeightedPoint.action_convolution, ComputableAnalysis.WeightedPoint.action_convolution_comm, ComputableAnalysis.WeightedPoint.action_convolution_assoc, ComputableAnalysis.FiniteProbabilityKernel.variance_convolution, ComputableAnalysis.FiniteProbabilityKernel.variance_convolutionPower, ComputableAnalysis.FiniteProbabilityKernel.convolveFunction_convolution}
\leanok
\end{theorem}

\section{A quantitative replacement step}

Write $\rho_3(K)=\sum_i p_i|x_i|^3$. Equal first and second moments imply
equal actions on every quadratic polynomial. This cancels the first three
terms in a local quadratic approximation; only the remainder contributes.

\begin{theorem}[Finite rescaled convolution comparison]
\label{thm:finite-convolution-replacement}
Let $K,L$ have equal first and second moments. Suppose rational $r,C$ and
quadratics $P_z$ satisfy
\[
 |h(z+rx)-P_z(x)|\le C|r|^3|x|^3
\]
on both finite supports for every rational $z$. Then
\[
 \left|A_{K^{*n}}(h(r\,\cdot))-A_{L^{*n}}(h(r\,\cdot))\right|
 \le nC|r|^3\bigl(\rho_3(K)+\rho_3(L)\bigr).
\]
\lean{ComputableAnalysis.FiniteProbabilityKernel.quadratic_action_eq, ComputableAnalysis.FiniteProbabilityKernel.cubic_replacement_bound, ComputableAnalysis.FiniteProbabilityKernel.convolutionPower_replacement_bound, ComputableAnalysis.FiniteProbabilityKernel.scaled_convolutionPower_cubic_bound}
\leanok
\end{theorem}

For $n=m^2$ and $r=1/m$, the coefficient is $C/m$. This theorem is finite
and rational. It assumes the stated remainder estimate, not differentiability,
and does not yet assert convergence to a Gaussian.

\section{Gaussian and central limit targets}

The existing Gaussian work supplies bounded quadrature comparisons,
reciprocal-square tail domination, and a valid full-line candidate with
width at most $16/(n+1)$. The candidate's compact integral witnesses and
exhaustion semantics still need a public comparison proof. Its identification
with $\sqrt\pi$ is also a target. The desired variants include, for $a>0$,
\[
 \int_{-\infty}^{\infty}e^{-ax^2+bx+c}\,dx
 \simeq\sqrt{\frac\pi a}\,e^{c+b^2/(4a)},
 \qquad G_s*G_t\simeq G_{s+t},
 \qquad G_t(x)=\frac{e^{-x^2/(4t)}}{\sqrt{4\pi t}}.
\]
These are not checked consequences of raw validity. They require justified
integrals, finite change-of-variables comparisons and tail bounds, with
arbitrary valid represented parameters and representation invariance.

For a nonnegative density $f$ of mass one, mean zero and variance one, put
$f_n(x)=\sqrt n\,f^{*n}(\sqrt n x)$. The first central limit target is
\[
 \int h(x)f_n(x)\,dx\longrightarrow
 \int h(x)\frac{e^{-x^2/2}}{\sqrt{2\pi}}\,dx
\]
for smooth bounded test functions. Pointwise convergence of densities is a
stronger, separate target. A supplied third absolute moment and cubic Taylor
remainder should yield an explicit rate through the checked replacement
theorem. The remaining work is continuous quadrature transport, Gaussian
stability and moments, and represented scaling. For the finite-variance
version, a second-moment tail modulus supplies the truncation schedule.
No measure-space or completed-real construction is required by this plan.
